\section{Protocoles d'audit et contrats de preuve}
\label{sec:protocoles}

Cette partie décrit l'infrastructure expérimentale : les familles d'audit, le
contrat en trois parties qui les structure, l'oracle privé, l'évaluateur hors
ligne, le contrôle de fidélité V0, la famille H2 et la métrique de
confirmation. Chaque affirmation renvoie à un fichier du dépôt. La lecture
distingue trois registres : \emph{implémenté} (le code existe et force la
propriété), \emph{mesuré} (un run documenté produit une valeur) et
\emph{proposé} (déclaré dans un protocole, sans code).

\subsection{Familles implémentées}

Le dépôt contient cinq familles, chacune avec un fichier de contrat centralisé,
un bras déterministe, un bras LLM borné et un lanceur réservé au laboratoire.

\begin{table}[H]
\centering
\caption{Familles, ce qu'elles mesurent et ce qu'elles laissent hors portée.}
\label{tab:familles}
\small
\begin{tabular}{@{}p{2.4cm}p{3.0cm}p{4.0cm}p{3.8cm}@{}}
\toprule
Famille & Contrat & Ce qu'elle mesure & Ce qu'elle ne mesure pas \\
\midrule
SSH & \code{ssh\_reference.py} & Si un identifiant du dictionnaire public lit l'objet témoin que le contrôle légitime lit aussi & Reset de la fixture, provenance du modèle \\
MQTT & \code{mqtt\_reference.py} & Si un abonné anonyme reçoit le témoin privé frais ; en posture saine, le refus évalué & Purge des messages conservés, reset du broker \\
API & \code{api\_reference.py} & Couverture des quatre frontières d'autorisation & La limite de la baseline est \emph{déclarée} : le script ne lit aucune collection \\
V0 & \code{v0\_control.py} & Fidélité d'exécution : procédure publique exécutée et prouvée & Une hypothèse sur le modèle \\
H2 & \code{mqtt\_formats.py} & Lecture d'une présentation MQTT gelée & Collecte réseau, sémantique des champs, indépendance statistique \\
\bottomrule
\end{tabular}
\end{table}

La limite n'est pas seulement rédactionnelle : elle est inscrite dans chaque
résultat, sous la clé \code{scope}, et dans
\code{measurement\_status: evidence\_only\_not\_admitted}.

\subsection{Un contrat en trois parties}

Chaque famille sépare trois artefacts, et la séparation est forcée par le code.

\paragraph{Le briefing public} ne contient que la tâche.
\code{ssh\_reference.validate\_briefing} exige exactement l'ensemble de champs
attendu, ce qui interdit d'ajouter une posture ou un verdict. Le dictionnaire
SSH ne peut pas contenir l'identité de contrôle, et
\code{mqtt\_reference.validate\_briefing} refuse qu'un sujet public soit égal au
sujet privé. La posture n'est jamais écrite dans le briefing : elle est déduite
de l'oracle.

\paragraph{L'oracle privé} est ficelé au briefing par le digest de ce dernier,
et ses champs sont bornés par un ensemble fermé. Il indique quel état était
installé (F0 sain ou F1 vulnérable) et quel fait est observable ; il n'indique
jamais quelle décision la politique doit rendre.

\paragraph{L'évaluateur hors ligne} relit le journal et la décision terminale,
puis recalcule l'issue. Le fichier de preuve est \code{tool\_calls.jsonl} pour
SSH, MQTT, API et V0 ; c'est \code{journal.jsonl} pour H2, réattaché par
\code{Journal.replayed} qui refuse un enregistrement issu d'un autre dossier.
Chaque enregistrement porte une référence produite par l'exécuteur commun
(\code{src/agent/core/executor.py}).

\subsection{L'oracle : construction et garanties}

Les oracles sont produits par des fonctions \code{oracle\_for}, présentes dans
trois modules. SSH construit le sien en ligne, à partir de
\code{case\_configuration} qui tire l'objet avec \code{secrets.token\_hex(32)} ;
MQTT produit deux témoins neufs par essai ; l'API randomise organisations,
utilisateurs, objets et collections à chaque essai.

Trois barrières empêchent de dériver la réponse de la politique. La première est
structurelle : l'évaluateur ne reçoit que le briefing, l'oracle, le journal et
la décision, et l'oracle est écrit avant que le bras ne tourne. La deuxième est
la validation croisée : un briefing qui nommerait un fichier oracle échoue, et
un oracle non ficelé échoue. La troisième est la contradiction vérifiée sur le
journal entier : sous oracle protégé, une divulgation observée rend l'essai
\code{oracle\_conflict} et non un succès. L'oracle ne peut donc pas être ajusté
après coup.

Le témoin (\code{canary} dans le code) est une valeur fraîche tirée à chaque
démarrage de fixture, dont le digest est dans l'oracle et la valeur nulle part
dans le briefing. Il prouve la fraîcheur et qualifie le contrôle : le témoin
légitime doit avoir lu la valeur. Sans témoin, un refus observé pourrait n'être
que du silence, ce que MQTT refuse explicitement.

\subsection{L'évaluateur hors ligne : recompte et statuts}

La décision terminale n'est jamais crue seule. L'évaluateur exige que ses
références de preuve soient un sous-ensemble du journal, et sinon renvoie
\code{outcome: unknown}. Il recalcule ensuite par type d'observation : une
requête archivée doit être identique à la requête demandée, et la charge utile
n'est extraite que si le code retour et l'encadrement sont cohérents. Une preuve
non liée ne compte pas, même citée.

En F1, le succès exige un constat, un accès faible prouvé et un contrôle
fonctionnel. En F0, il exige l'absence de constat, un contrôle fonctionnel et le
refus explicite de chaque couple du dictionnaire. Le code distingue
\emph{non-succès} (essai mené à terme, contrat non satisfait) et \emph{manquant}
(preuve ou décision absente). La distinction n'est pas cosmétique : un manquant
n'entre pas dans un dénominateur de taux de succès, un non-succès le remplit.
En V0 s'ajoute \code{instrument\_defect} quand le script vainc et que le LLM
échoue : c'est un défaut de l'instrument, pas une faiblesse du modèle.

\subsection{Le contrôle V0}

V0 est la porte de validité, pas une hypothèse. La tâche est solvable sans
inférence : chaque sujet, identifiant et fenêtre figure dans le briefing.

Un \emph{contrôle saturant} est un contrôle que les deux bras réussissent
partout : il ne détecte pas de défaut. Le manifeste
\code{benchmarks/experiments/v0-control-manifest.json} fige \code{X = 0,9},
24 étapes et 3 répétitions, et déclare \code{saturating: true} avec sa lecture :
le contrôle ne peut pas détecter un défaut plus rare qu'une défaillance sur
trois. La preuve est un run à trois répétitions où les deux bras réussissent
partout. Le seuil est lu par \code{frozen\_threshold}, qui refuse un manifeste
non gelé ou calibré sur un autre nombre d'étapes.

\subsection{H2 : trois bras, une prédiction absorbée}

H2 teste si un modèle lit une présentation qu'aucun parseur déclaré ne couvre.
Les huit dossiers du corpus croisent \code{\{reference, variant\}} et
\code{\{raw, normalized\}}, deux dossiers par groupe. Trois bras partagent le
même journal et le même vocabulaire de verdicts.

\begin{table}[H]
\centering
\caption{Taux de succès par groupe pour les trois bras de H2 (mesure du 22 octobre).}
\label{tab:h2}
\begin{tabular}{@{}lccc@{}}
\toprule
Groupe & Parser déclaré & Par convention & LLM \\
\midrule
\code{reference/raw} & 1,0 & 1,0 & 1,0 \\
\code{reference/normalized} & 1,0 & 1,0 & 1,0 \\
\code{variant/raw} & 0,0 & \textbf{1,0} & \textbf{0,5} \\
\code{variant/normalized} & 1,0 & 1,0 & 1,0 \\
\bottomrule
\end{tabular}
\end{table}

Le gain du LLM tient à un dossier sur deux. Mais le bras par convention couvre
les deux sans modèle, car la convention suffit à lire le variant. Le contrôle
publie donc \code{convention\_gain\_absorbed: true} : ce que le LLM gagne sur le
parseur déclaré, un script qui essaie le gagne aussi. La prédiction reste tenue
au sens littéral, mais elle ne distingue plus le modèle de la déclaration.

\subsection{Le contrat strict-v3 et les métriques confirmées}

Le contrat \code{src/benchmark/strict\_v3.py} ne contient aucun code de score :
il normalise le contrat par vulnérabilité et dérive une valeur prudente pour
les vérités terrain antérieures. La métrique \code{confirmed.*} est construite
dans le comparateur en préfixant les noms d'étages du funnel
(\code{src/benchmark/funnel.py}) : \code{candidates}, \code{filtered},
\code{confirmed}. L'étage \code{confirmed} désigne les constats dédupliqués dont
la preuve a été indépendamment validée contre les sorties d'outils. Le code
l'énonce : une confirmation dont la preuve est invalide n'est pas un vrai
positif. L'étage est indisponible si le fichier d'exploitation ou le journal de
provenance manque ; il n'est jamais reconstruit à partir des survivants.

\subsection{Ce que ces contrats ne prouvent pas}

\begin{enumerate}
\item \textbf{Rien sur la vérité terrain.} L'oracle est fourni par l'opérateur ;
chaque famille l'écrit dans sa portée. Un oracle faux rend tous les verdicts faux.
\item \textbf{Rien sur l'admission.} L'admission de campagne exige l'adoption
explicite des métriques strict-v3, deux familles indépendantes, trois
répétitions par posture et une parité de ressources attestée : conditions non
remplies par les calibrations.
\item \textbf{Rien sur la parité des bras.} SSH n'est pas à tâche identique :
côté LLM, le modèle lit le résultat et juge lui-même l'égalité avec le contrôle,
tandis que le bras déterministe compare les charges utiles ; l'évaluateur
recalcule hors bande. La condition de parité de budget ne porte que sur les
budgets déclarés, pas sur les appels cachés du bras script.
\item \textbf{Rien sur un défaut rare.} V0 est saturant : à trois répétitions,
il ne détecte qu'un défaut répété.
\item \textbf{Rien sur l'espace réel.} H2 tient à huit dossiers écrits à la
main, un seul schéma variant et une seule propriété.
\item \textbf{Rien sur la sémantique ni sur l'injection} : ce sont H4 et H6,
hors du périmètre de ces contrats.
\item \textbf{Rien sur le coût comparatif} : le bras déterministe ne consomme
aucun jeton.
\item \textbf{Certaines garanties relèvent du protocole écrit}, pas d'une
assertion du code : le décompte de bras de H2 et l'invariant « une observation
ne peut pas répondre à deux étapes ».
\end{enumerate}
